\chapter{Language, Response Format, and Conjunction Judgments}
\label{sec:p19-conjunction-language}

The probability of an intersection cannot exceed the probability of
either constituent under the same information. Applying that benchmark
to a verbal response also requires identifying what the respondent
understood and what the response measured.
Chapter~\ref{sec:p19-conjunction-judgments} supplies the mathematical
benchmark and the original 1983 experiments. This chapter presents
Hertwig and Gigerenzer's 1999 investigation of language and response
format \cite{HertwigGigerenzer1999Conjunction}.
Its four numbered studies, additional comparisons, and discussion of
earlier work are kept distinct.

The paper proposes that ordinary-language interpretation can explain
some responses classified as conjunction errors. Its data concern
particular elicitation procedures, chiefly with German-speaking
university students. Its broader defense of social rationality is an
argument about interpreting those findings; it is not an experimentally
established diagnosis of every conjunction response.

\section{Separate the benchmark, the question, and the response}
\label{sec:p19-conjunction-language-targets}

\begin{hbdefinition}{def:p19-conjunction-language-record}{An interpretation-sensitive response record}
For each task record the described person or population, the information
given, the alternatives, the requested quantity, the response format,
the order of questions, and the interpretation checks.
In the Linda comparisons let \(T\) denote bank-teller work and \(F\)
feminist activity. A probability comparison targets
\(P(T\cap F\mid E)\) and \(P(T\mid E)\), with the same information
\(E\), \(P(E)>0\). A frequency comparison targets the numbers satisfying
the two descriptions within the same collection of cases.

Probability ranking, numerical probability estimation, numerical
frequency estimation, typicality rating, and a verbal explanation are
different response formats. A participant's subsequent paraphrase is
another observation. It should be stored with its timing and coding rule,
without assuming that it is an infallible report of the process that
produced the preceding judgment.

The 1999 paper calls a strict reversal of the intended inclusion order
a \emph{conjunction violation}. Its principal empirical question is
whether such an ordering was produced while answering the intended
mathematical question. A probability error and a disagreement about the
question's meaning require different explanations.
\index{conjunction fallacy!language and response format}
\index{semantic inference}\index{response format}
\end{hbdefinition}

The source uses \emph{semantic inference} for inferring a word's
intended meaning from its surrounding content and conversational
context. Its account draws on expectations of relevance and adequate
informativeness. On that account, an elaborate personality description
encourages respondents to use the description, and an immediately
repeated request is normally expected to ask for something different.
The authors propose that these expectations can favor readings of
\emph{probability} concerned with plausibility, resemblance, or
believability rather than a numerical probability measure.

These are proposed psychological connections. A word's dictionary
alternatives do not prove which meaning a particular person used.
Conversely, a mathematical probability interpretation does not require
discarding all information in the description. Information can affect
both numerical probabilities and the ordering of the two nonnested
constituents even though it cannot reverse inclusion.

\begin{hbexample}{def:p19-conjunction-profile-relevance}{Relevant information with unchanged inclusion constraints}
Partition the sample space into two stipulated information conditions
\(E_1,E_2\), each with probability \(1/2\). Within each assign the following
conditional masses to the four cells:
\[
\begin{array}{c|rrrr}
 &T\cap F&T^c\cap F&T\cap F^c&T^c\cap F^c\\ \hline
 E_1&.30&.50&.10&.10\\
 E_2&.10&.10&.60&.20
\end{array}
\]
Multiplying each row by \(1/2\) defines eight nonnegative atom masses
whose sum is one. Thus this is one coherent probability model, not
two inconsistent assignments.
Within \(E_1\), the probabilities of \(F,T,T\cap F\) are
\(.80,.40,.30\); within \(E_2\), those of \(T,F,T\cap F\)
are \(.70,.20,.10\). The condition changes the values and the
ordering of \(T\) and \(F\), while the conjunction remains no more
probable than either.

This constructed example makes no claim about Linda or a participant.
It shows why the inference from an expectation of relevance to a
nonprobabilistic interpretation is a substantive psychological
hypothesis. The 1999 authors acknowledge the \(T\)-versus-\(F\)
point and reply that relevance can vary by degree
\cite[p.~297]{HertwigGigerenzer1999Conjunction}.
That reply belongs to their proposed interpretation mechanism.
\end{hbexample}

\section{The three predictions}
\label{sec:p19-conjunction-language-predictions}

The source derives three empirical predictions from its semantic account
\cite[pp.~278--279]{HertwigGigerenzer1999Conjunction}.

\begin{enumerate}
\item Ordinary probability wording should often elicit interpretations
outside the intended mathematical comparison, accompanied by many
inclusion reversals. Asking for typicality first should reduce reversals
in the subsequent probability judgment: the description has already
served an explicit purpose, and the second question invites a different
assessment.
\item Frequency wording should elicit interpretations involving numbers
of cases and fewer reversals. This reduction should persist when the
task omits an explicit population size or first attaches the descriptive
features to one person rather than a group.
\item Replacing probability wording by believability wording should
produce similar rankings if believability is already a plausible
interpretation of the original question.
\end{enumerate}

The competing account discussed in the article attributes frequency
effects to cues that make set inclusion easier to represent.
The 1999 experiments operationalize two such cues: attaching features
to a reference class and stating that class's numerical size.
The experiments concern these specified manipulations; they do not
exhaust every possible way of making inclusion salient.
Response mode is an additional factor: ranking descriptions and
estimating numbers can elicit different strategies.

\section{Study 1: probability paraphrases after a Linda ranking}
\label{sec:p19-conjunction-language-study-one}

Eighteen University of Munich students, recruited by advertisement
across disciplines, participated individually. None reported having
encountered the Linda problem. They first ranked \(T,F,T\cap F\)
by probability, with alternative order counterbalanced. They then
explained the task to an imagined person whose German did not include
the word for probability. Their oral paraphrases were recorded,
replayed, and refined into instructions. Finally, they selected one
or more meanings from thirteen alternatives, with room to add another
and with two randomized list orders
\cite[pp.~279--283]{HertwigGigerenzer1999Conjunction}.

The offered alternatives concerned logicality, certainty, frequency,
typicality, credibility, plausibility, tenability, evidence,
conceivability, possibility, predictability, similarity, and
reasonableness. Most came from an earlier pilot; the list was not
balanced between the authors' two coding classes.
The authors classified paraphrases as mathematical when they judged
them to express one of their selected classical, frequency, logical,
or subjective interpretations. Thus their category includes responses
such as certainty and logicality. This is the study's coding convention;
the labels alone do not demonstrate a coherent probability assignment.

Fifteen of the eighteen participants, reported as 83\%, reversed an
intended inclusion comparison. The same eighteen produced 39
paraphrases, covering eighteen coded meanings. Seven paraphrases,
18\% of responses, were coded mathematical; 32 were coded otherwise.
The original figure's counts are:

\begin{description}
\item[Paraphrases coded mathematical.]
Certainty occurred twice; frequency, expected value, logicality,
percentage, and randomness once each.
\item[Other paraphrases.]
Possibility occurred nine times, correspondence six, applicability
and conceivability four each, and believability twice.
Credibility, importance, obviousness, plausibility, reality,
similarity, and comprehensibility occurred once each.
\end{description}

The 39 is a count of responses, not independent participants.
Two of the three participants without a reversal produced at least
one mathematical paraphrase. Twelve of the fifteen with a reversal
produced none. Correspondence was a heterogeneous category involving
relations between the profile, worldview, and activities.
Combining it with similarity yields seven of the 32 responses in the
other category, the source's 22\% comparison; this regrouping is not
a direct observation of a representativeness computation.

The subsequent multiple-choice task produced 64 selections:
eight mathematical and 56 other selections, the latter reported as
88\%. Certainty was selected five times, logicality twice, and
frequency once. Among the other alternatives, counts were possibility
11, conceivability 9, plausibility and reasonableness 8 each,
typicality 6, predictability 4, credibility, evidence, and similarity
3 each, and tenability 1. Dividing by the numbers of offered
alternatives gives the authors' averages of about 2.7 selections per
mathematical alternative and 5.6 per other alternative.
That comparison addresses the list imbalance; it does not turn the
selected meanings into mutually exclusive participant groups.

The later discussion reports an additional interpretation question
to the Study 1 participants: six of seventeen respondents, 35\%,
understood the unqualified teller option as excluding feminist
activity; one person did not answer. All six also supplied at least
one nonmathematical probability paraphrase.
The two proposed interpretation changes can therefore coexist in
this sample. They are not independent explanations whose percentages
can simply be added.

\section{Study 2: typicality first and thinking aloud}
\label{sec:p19-conjunction-language-study-two}

Ninety University of Salzburg students, recruited across disciplines
and unfamiliar with the Linda problem, were randomly assigned to
four conditions. The two typicality-first conditions introduced
both questions together. Participants rated the fit of Linda to
each of the three categories on a seven-point scale and then ranked
the three descriptions by probability. The typicality instructions
distinguished category representativeness from liking the category.
The other conditions asked only for the probability ranking.

Thinking aloud was crossed with this manipulation. Participants
assigned to it practiced on two problems; their protocols were
recorded and transcribed. The probability-only thinking-aloud group,
the typicality-first silent group, and the typicality-first
thinking-aloud group were tested individually.
Table~\ref{tab:p19-conjunction-typicality-first} gives the complete
reported allocation and critical outcomes
\cite[pp.~283--285]{HertwigGigerenzer1999Conjunction}.

\begin{table}[htbp]
\centering
\begin{tabular}{p{.52\textwidth}p{.24\textwidth}p{.14\textwidth}}
\hline
Condition & Reversals / respondents & Reported percent\\ \hline
Probability ranking only, silent &21/24&88\%\\
Probability ranking only, thinking aloud &18/22&82\%\\
Typicality then probability, silent &13/22&59\%\\
Typicality then probability, thinking aloud &6/22&27\%\\ \hline
\end{tabular}
\caption{Study 2: reversals of the intended conjunction ordering.
The four randomized conditions together contain 90 participants.}
\label{tab:p19-conjunction-typicality-first}
\end{table}

The source reports reductions of 29 percentage points without thinking
aloud and 55 points with it, with effect-size values \(f=.32\)
and \(f=.55\). These are the paper's rounded summaries.
It does not explain why the reduction was larger with the combined
procedure. The 88\%-versus-82\% comparison alone does not establish
equivalence of silent and thinking-aloud elicitation.

The protocol categories also matter. In the typicality-first
thinking-aloud condition, sixteen participants ranked the conjunction
last. Twelve of them supplied one of two kinds of justification:
five invoked class inclusion and seven considered feminist activity
and teller work incompatible, making the intersection empty.
Both routes imply the observed ordering, but the second introduces
a substantive incompatibility premise that is not supplied by set
inclusion alone.

In the probability-only thinking-aloud condition, twelve of the
eighteen who ranked teller work last argued that some profile
feature was incompatible with being a teller.
Of the four who ranked the conjunction last, one invoked inclusion
and two treated the intersection as empty.
These are conditional counts within the displayed response groups.
They are not 12 additional participants in either condition.
The protocols support a difference in articulated reasons; they
do not by themselves identify a single causal mental process.

\section{Study 3: frequency paraphrases}
\label{sec:p19-conjunction-language-study-three}

Twenty further Munich students, recruited across disciplines and
unfamiliar with the task, were tested individually.
They estimated the numbers satisfying \(T,F,T\cap F\) in a
hypothetical group of 200 women sharing the adapted profile, with
alternative order counterbalanced. The source group description
uses an average age of 30 and the education, social-concern, and
demonstration features; it is not literally the single 31-year-old
Linda wording. Participants then performed the imagined-listener
paraphrase task for frequency
\cite[pp.~287--288, Exhibit~3]{HertwigGigerenzer1999Conjunction}.

None of the twenty reversed inclusion in the numerical estimates.
They produced 55 paraphrases, spanning eleven meanings.
The authors coded 54 as referring to numerosity and one, possibility,
as not doing so. The counts were: how often 13, number 12,
percentage 11, proportion 6, quantity 5, probability 3, and one
each for expectancy value, favorable/possible events, regularity,
extent, and possibility. Some categories were closely related.
The authors note that how often resembles the how-many phrasing
already supplied in the task.

The zero observed reversals is a property of this twenty-person
sample. It does not establish a zero population rate.
The 54 of 55 concerns utterances coded under this task's rule.
It is neither 54 people nor a measure of their numerical accuracy:
the hypothetical group had no observed outcome counts against
which their estimates could be calibrated.
Study 1 and Study 3 used different samples and different judgment
formats; their contrast combines wording with ranking versus
numerical estimation.

\section{Study 4: separating two cues from the requested judgment}
\label{sec:p19-conjunction-language-study-four}

Ninety-eight Salzburg students, recruited across disciplines and
unfamiliar with the Linda problem, were randomly allocated to four
versions and tested in groups of at most six. Every version included
both constituents and their conjunction, with alternative order
counterbalanced. The versions were:

\begin{enumerate}
\item A group-first frequency question: the adapted features describe
200 women, and the respondent estimates the number satisfying each
description.
\item A person-first frequency question: Linda is described first,
then a hypothetical group of 200 matching women is introduced,
and the same numbers are requested.
\item A person-first frequency question with no stated group size:
the respondent imagines matching women and fills in both a count
and its reference total.
\item A probability-estimation question with both proposed cues:
200 women share the features; one is selected at random and named
Linda; numerical probabilities are requested for her descriptions.
\end{enumerate}

The fourth version retains a numerical estimation response and both
reference-class cues while changing the question from numbers of
women to a randomly selected woman's probabilities.
The complete Linda results are in
Table~\ref{tab:p19-conjunction-language-cues}
\cite[pp.~289--292, Exhibit~3]{HertwigGigerenzer1999Conjunction}.

\begin{table}[htbp]
\centering
\begin{tabular}{p{.52\textwidth}p{.24\textwidth}p{.14\textwidth}}
\hline
Question version & Reversals / completed responses & Percent\\ \hline
Group-first frequency, stated total &3/23&13\%\\
Person-first frequency, stated total &5/25&20\%\\
Person-first frequency, unstated total &4/25&16\%\\
Probability estimate with both cues &16/24&67\%\\ \hline
\end{tabular}
\caption{Study 4: the 97 completed Linda responses.
One participant in the group-first frequency condition did not
complete the problem; the recruited allocation was 98.}
\label{tab:p19-conjunction-language-cues}
\end{table}

Removing either or both of these two cues did not reproduce the
high reversal rate of the probability-estimation version in these
samples. The source reports the latter as 54 percentage points
above the group-first frequency version, \(f=.55\), and 47 and
51 points above the other frequency versions, \(f=.47,.52\).
It also compares these results with the probability-ranking
condition of Study 2: the probability-estimation version is 21
points lower, \(f=.25\), while the two person-first frequency
versions are 68 and 72 points lower, \(f=.68,.72\).
Those comparisons with Study 2 cross studies; they are not
additional randomized cells in Study 4.

The authors take the pattern to support semantic interpretation
and a separate response-mode contribution. Its scope is narrower
than a claim that reference-class information never helps.
The four conditions are not a full factorial manipulation of every
wording, cue, response mode, and task. Nor do the observed 13\%,
20\%, and 16\% establish exact equality of the three underlying
response probabilities.

\subsection{The additional health and tennis comparisons}

The source also reports questions about a man's heart-attack history
and age, and a tennis outcome involving Boris Becker, updating
the earlier Borg scenario. Its procedure note says that these
additional problems followed Linda for participants receiving them.
In probability versions the reported reversal counts are 15/24,
63\%, for health, and 14/24, 58\%, for tennis.
The comparison Linda probability result is the 21/24, 88\%,
Study 2 condition.
Frequency-version counts are 3/24, 13\%, for health; 2/23, 9\%,
for tennis; and 3/23, 13\%, for Linda in Study 4.

The reported differences from the Linda probability result are
25 and 30 percentage points, with \(f=.29,.33\).
The article does not print a complete participant-by-condition
allocation for every additional item or the complete new
Becker stimulus. These counts should not be pooled as independent
new samples or used to reconstruct missing wording.
The authors propose that different task content can direct
interpretation differently; the observed pattern does not isolate
content from every other procedural difference
\cite[p.~299 and p.~289 n.~6]{HertwigGigerenzer1999Conjunction}.

\section{An additional word-frequency comparison}
\label{sec:p19-conjunction-language-words}

The 1999 authors also investigated the original seven-letter-word
contrast. Three groups of 31 received questions about approximately
2,000 words in four novel pages. One group estimated both the
number ending in \emph{ing} and the number with \emph{n} in
position six. Each of the other groups estimated one category.
The paper gives these group sizes but does not identify the
recruitment population for this additional comparison in that passage.

In the separate groups, mean estimates were 425 for the narrower
suffix class and 212 for the broader sixth-position class;
the medians were 300 and 50. In the group answering both
questions, the respective means were 282 and 457 and the
medians 150 and 250. Eight of the 31 individuals in the latter
condition, 26\%, reversed inclusion
\cite[pp.~295--296]{HertwigGigerenzer1999Conjunction}.

These values are reported estimates, not actual counts in a
specified novel. Their absolute accuracy is not established.
The common-group ordering and the eight individual reversals
can coexist: an average respecting inclusion does not imply
that every participant respects it.
The separate-group reversal does not prove that each individual
would violate inclusion if asked both questions; the precise
aggregation distinction is proved in
\cref{thm:p19-conjunction-aggregation}.

\section{Competing interpretations and the source's replies}
\label{sec:p19-conjunction-language-objections}

\paragraph{Believability and the evidence actually available.}
The third prediction was assessed using Macdonald and Gilhooly's
earlier published summaries, rather than a fifth new experiment.
As reported by the 1999 article, mean ranks for \(T,F,T\cap F\)
were \(5.9,2.0,4.6\) under probability wording and
\(5.7,1.9,4.8\) under believability wording.
The earlier paper reportedly gave 89\% reversals for probability
but no corresponding believability reversal rate.
Similar mean ranks therefore support only that summary comparison;
they cannot establish an unreported individual reversal percentage
or statistical equivalence
\cite[p.~292]{HertwigGigerenzer1999Conjunction}.
The earlier original remains a separate source presentation.

\paragraph{Interpreting an unqualified option exclusively.}
Another explanation changes the event, reading \(T\) as
\(T\cap F^c\). Then \(T\cap F\) and \(T\cap F^c\)
are disjoint alternatives, and inclusion does not determine
their ordering. The concrete probability construction in
\cref{thm:p19-conjunction-interpretations} demonstrates this
possibility.

The 1999 discussion reports mixed outcomes from earlier attempts
to block such an interpretation: explicitly inclusive wording
reduced the original Linda result to 57\%; a separate
Messer--Griggs formulation reportedly changed 77\% to 56\%;
and adding an explicit nonfeminist alternative in
Morier--Borgida left Linda near its baseline, 77\% versus 80\%,
while reducing Bill from 77\% to 49\%.
It also contrasts Agnoli--Krantz's unsuccessful predicted
rephrasing pattern with Macdonald--Gilhooly's reduction from
75\% to 21\%, and Politzer--Noveck's 29\% on an exam
partition compared with 77\% on Linda.
These are the 1999 article's reports of distinct studies,
with different wording and designs
\cite[pp.~297--298]{HertwigGigerenzer1999Conjunction}.
They do not constitute one pooled estimate.

The same discussion reports Dulany--Hilton's retrospective
interpretation check: among 61 participants, 32 reversed
inclusion and sixteen of those selected an explicitly inclusive
reading. Their narrower diagnostic criterion therefore counted
16/61, about 26\%, as established violations.
The criterion and its timing remain part of that source's
method; they are not an automatic correction factor for
every later conjunction experiment. The complete earlier
originals and their alternative explanations require their
own treatments.

\paragraph{Betting and a different event partition.}
The actual original twenty-roll betting instruction omits the
word probability. Nevertheless, 65\% and 62\% selected
\texttt{GRGRRR} over its contained shorter pattern.
The 1999 discussion proposes another possible reading:
counting red and green outcomes without their order.
For fixed blocks of five or six rolls, respectively, this
gives the following binomial calculations:
\[
\begin{aligned}
 P(\text{four reds in five})&={5\choose4}(1/3)^4(2/3)
                                      =10/243,\\
 P(\text{four reds in six})&={6\choose4}(1/3)^4(2/3)^2
                                      =20/243,\\
 P(\text{five reds in six})&={6\choose5}(1/3)^5(2/3)
                                      =4/243.
\end{aligned}
\]
Thus the source's approximately \(.041,.082,.016\)
comparison can hold on that different target.
Each factor counts placements of the red outcomes in its
specified block, and each term uses that block's length.
These are not probabilities of the printed strings appearing
somewhere in twenty rolls. The latter probabilities and their
strict dominance relation are derived in
\cref{thm:p19-conjunction-bets}.
An alternative calculation shows logical feasibility of a
reading; evidence that the original participants adopted
that reading is a separate requirement.

\paragraph{Probability, information, and a requested answer.}
The paper's discussion of Bar-Hillel--Neter's geographical
questions distinguishes a probable event from an informative
answer at the requested level of specificity.
A country answer and a continent answer can differ in
communicative usefulness even when the country event is a
subset of the continent event. The reported examples and
betting variants are earlier-source reports, not new 1999
experiments. A desire for specificity does not reverse
the mathematical inclusion of the events under fixed payouts;
it changes the objective used to evaluate an answer.

\paragraph{Explicit statistical instructions and other cues.}
The article discusses Donovan--Epstein's instruction to treat
the problems as statistical exercises. It reports an aggregate
reversal level roughly thirty points below a historical Linda
median, while acknowledging those authors' contrary conclusion
about conversational explanations.
That historical-median comparison is not the within-study
causal effect of the instruction.

The 1999 authors also propose further cues: numerical base rates
or likelihoods, an explicit sampling mechanism, a requirement
to calculate, and a familiar problem class such as balls and
urns. These suggestions extend the hypothesis beyond the two
conversational maxims and the two cues experimentally
manipulated here. The paper explicitly acknowledges the lack
of a complete process model of semantic inference
\cite[pp.~293--298]{HertwigGigerenzer1999Conjunction}.

\section{What a frequency comparison does and does not isolate}
\label{sec:p19-conjunction-language-comparison}

The article's appendix makes a useful source comparison.
The original health question changes a randomly selected
person's probability comparison to counts in a group of 100.
Fiedler's version first describes Linda and then asks about
100 people like her. Reeves--Lockhart's example supplies
attendance frequencies for several named people but still
asks for rankings of probabilities that people are present
at a bar. Frequency information supplied as evidence and a
frequency judgment requested as a response are different
manipulations.

The article accordingly cautions against treating earlier
studies carrying a frequency label as the same experiment.
It reports sizeable variation among their outcomes, and
notes that Jones and colleagues did not supply enough
wording to settle whether the relevant output was a
frequency or percentage judgment. Its response-mode
discussion draws on Hertwig--Chase and Erdfelder and
colleagues; those are separately attributed prior studies,
not additional participants in Studies 1--4.
The appendix's descriptions preserve the selected earlier
versions, but do not replace acquiring and reading each
complete original \cite[pp.~285--287,304--305]{HertwigGigerenzer1999Conjunction}.

Frequency representations also play different roles in other
tasks. Natural-frequency Bayesian calculation can reduce the
number of arithmetic steps, whereas a conjunction question
may principally change interpretation or expose inclusion.
Calibration across outcomes and coherence among simultaneous
judgments answer different questions. The separate treatments
in \cref{sec:p19-calibration-extreme-confidence,sec:p19-natural-frequencies-bayesian-reasoning}
should be used
with their own task definitions.
The 1999 authors explicitly reject the conclusion that
frequency judgments are always accurate.

\section{Social context and the limits of the inference}
\label{sec:p19-conjunction-language-limits}

The source closes with the possibility that social purposes
make a choice intelligible even when it violates a condition
on an abstract choice function.
For precision, let \(C(M)\) be a nonempty set of chosen
alternatives from a finite nonempty menu \(M\).
Contraction consistency requires
\[
       x\in C(M),\quad x\in A\subseteq M
                           \quad\Longrightarrow\quad x\in C(A).
\]
In the source's dessert example, let \(x\) mean declining dessert,
\(y\) taking the eclair, and \(z\) taking a different pastry.
Choosing nothing from
\(\{x,y\}\) to leave the last pastry for someone else,
but choosing the pastry \(y\) from \(\{x,y,z\}\), gives
\(C(\{x,y\})=\{x\}\) and \(C(\{x,y,z\})=\{y\}\).
Since \(y\) belongs to both menus, these choices violate
the displayed condition.
The social consequence of taking the pastry changes with
what remains for another guest. A model of the person's
full consequences can therefore need more than the bare
food labels.

This example illustrates the importance of specifying
alternatives and goals; it is not evidence for the
semantic mechanism proposed for Linda.
The source invokes Sen's discussion of choice and Wundt's
distinction between normative logic and psychological
process. Those intellectual connections are identified
without treating their cited works as newly acquired
primary presentations.

For the reported experiments, the principal limits are
concrete. The recruited samples were small and restricted;
retrospective paraphrases and think-aloud instructions can
alter or incompletely reveal processing; coding a word as
mathematical does not demonstrate numerical coherence;
and numerical coherence does not establish factual accuracy.
The paper does not report confidence intervals for the
displayed proportions or a coder-agreement coefficient
for the paraphrase classification. Its effect-size
summaries should not be expanded into unreported analyses.
Comparisons across studies and previously published
conditions retain their different allocations and measures.

The 1983 evidence and the 1999 findings support a reference
practice that records task meaning and response procedure
alongside the inclusion benchmark. They do not support
either diagnosing global irrationality from one reversal
or declaring every reversal a successful semantic inference.
Later direct tests, response models, and replication studies
remain separately identified sources of evidence.

\begin{hbmethod}{def:p19-conjunction-language-method}{Assessing a proposed interpretation explanation}
Specify the intended event comparison and the exact question.
Write the alternative interpretation as a different event,
quantity, or objective; calculate what it predicts before
using it to explain a response.

Identify the evidence for that interpretation: an elicited
paraphrase, a choice among meanings, a concurrent explanation,
or a manipulation with predicted effects. Record when that
evidence was obtained and whether the check itself changed
the task. Preserve both participant counts and response
counts when one person supplies several utterances.

Separate wording, response mode, reference-class information,
question order, and incentives. Use the assigned conditions
for a within-study comparison and mark comparisons with
earlier samples as such. A small observed difference does
not establish equivalence; no observed errors in a finite
sample does not establish their impossibility.

Report the strongest supported conclusion. Evidence that
wording changes responses supports a formulation effect.
Evidence associating a paraphrase with a response supports
that association. Identifying a causal processing mechanism
requires an argument connecting the observations to that
mechanism and distinguishing alternatives.
When performing a probability calculation, retain the
common-event-space inclusion constraint. When studying
behavior, retain the original response and its interpretation
evidence rather than silently recoding it to the preferred
explanation.
\end{hbmethod}
